\section{Evidence map and artifact identities}
\label{app:artifacts}
The full identities are stored in accompanying JSON receipts and manifests.
The 12 November source-intent diagnostic is bound to \code{897c121} and
\path{docs/evidence/source-intent-reference-2026-11-12.json}; its initial and
corrective raw receipts remain in distinct run folders. The r10 viewer is
instead bound to \code{663b7e2} by
\path{docs/evidence/pdf-navigation-r10-artifacts-2026-11-12.json}. Newer core
draft/preamble APIs are not included in that earlier binary. Engineering checks
are separately recorded in
\path{docs/architecture/source-draft-validation-2026-11-12.json}.
Table~\ref{tab:artifact-map} maps the principal stages. These records are local
development artifacts; no public release or submission is implied.

\begin{table}[htbp]
\small\begin{tabularx}{\linewidth}{@{}>{\raggedright\arraybackslash}p{.29\linewidth}Y@{}}
\toprule
Record & Purpose / authority\\
\midrule
\code{c26a238} & Inspected audit baseline and earlier observed three-workflow CI\\
\code{5e59e77} & Historical r4 source and ordinary/diagnostic build evidence\\
\code{c2ef008}, \code{391332d} & Later failure handling and interaction completion; resumed tests recorded separately\\
\code{d267118}, \code{a120086} & Replacement verifier source and exact r5 artifact/validation checkpoint\\
\code{cb3cd11} & Opt-in native candidate-boundary harness source\\
\code{b5edb9d} & Corrected reference/health code, before its new run\\
\path{research/evidence-catalog.json} & Receipt fingerprints; reading a receipt is not rerunning its experiment\\
\path{research/evaluation/freeze.json} & Deliberately unset final-evaluation version/manifests\\
\bottomrule
\end{tabularx}
\caption{Code/evidence identities serve different roles.}
\label{tab:artifact-map}
\end{table}

The principal dated receipts under \path{docs/evidence/} are
\path{setup-timing-local-2026-11-05.json},
\path{native-reference-development-2026-11-05.json},
\path{pdfium-attestation-verification-2026-11-05.json},
\path{research-environment-2026-11-05.json}, and the separate report-draft
receipts. Earlier installed/native/engine results retain their own versions.
The inventory identifies both successful repeats and inconclusive/failed
attempts rather than replacing their meaning with the current status.

\subsection{Exact executable and verifier fingerprints}
The following SHA-256 values distinguish the actual Windows artifacts. The
debug native harness is not a tested installation of the ordinary GUI build.

\textbf{Ordinary r4 executable:}\\
\code{a9e7dfd5e523e5025ccc6f98da7ece7ad16ae254b11a4ae8da7b7135a4dc4ee3}

\textbf{Ordinary r4 NSIS setup:}\\
\code{6073346d7c65e903a4ba8c5a89ac26ab4990f891340bfb802be4d1a6c8396112}

\textbf{Diagnostic r4 executable:}\\
\code{2e4173437be96914a92747d087ce448423e8a56f1a6b94adbaf714e7dfaecce7}

\textbf{Native development harness, both attempts:}\\
\code{fcb46dae8fcdbc86e415209cc9f9caf62892ff4315009658185468768422c2a3}

\textbf{Pinned pypdfium2 Windows wheel:}\\
\code{149fd5c6397b8df8bf7911a93506eff0be874f877afe7ac936cf5d37d21a6a06}

\textbf{Packaged PDFium DLL:}\\
\code{69f1e860f6742fbaafcc4db80e785eec9a4ee0da36ea6d1dc0867ea0f18baea6}

\textbf{Packaged Rust inspection worker:}\\
\code{4c989986a8e1b89cac4ba4bf25043b9719be45b5ebead393d1d81ba027613840}

\subsection{Raw run locations and retention}
Pilot A is under \path{artifacts/oracle-feasibility-20261104-first/}.
The native-entry first attempt, corrected attempt and posthoc diagnosis are
under \path{artifacts/native-reference-development-20261105-first/},
\path{artifacts/native-reference-development-20261105-reference-v2/}, and
\path{artifacts/native-reference-diagnosis-20261105/}. The first native
protocol fingerprint is:

\code{8cd000d559011a5e14e7accb5dc1ab67967763dd936cb648ea079b625d5d6a84}

Archives retain synthetic input/source snapshots, PDFs, comparison images/XML,
logs, outcomes and compiler build manifests. Complete local caches remain in
the run directories. Portable evidence archives explicitly exclude executable/
DLL binaries, standalone font files, live locks and recovery internals where
those are not necessary to interpret the result. Their manifests enumerate
what is actually included; the archive is not described as a full machine image.

The partial environment and provenance checks are under
\path{artifacts/research-environment-20261105/} and
\path{artifacts/attestation-review-20261105/}. The latter keeps tool-install
hashes/versions, public provenance, verification logs, negative-control
metadata and a partial post-verification trust-cache snapshot. The copied
changed wheel was never installed or loaded, and the original runtime pin
remained unchanged.
